\cotefile{1}
\scanpage{93}

\uline{Th.~4} (\uline{troisième th. de convexité dans les espaces de Hilbert})

Soient $E,F$ deux espaces de Hilbert, $u,v$ deux opérateurs compacts de $E$ dans $F$, soit enfin $\varphi\inserted{(s_1,\ldots s_n)}$ une fonction \deleted{convexe symétrique}\inserted{convexe symétrique croissante \deleted{et \unsure{épist.}}} de $n$ arguments $s_i\geq 0$, \deleted{telle que \unsure{$f$ l'ét.|$f(e^{t_1},$}, \unsure{étant|$e^{t_n})$} soit fonction convexe de $(h)$ $\in R^n$. Alors,}\inserted{Alors}

\[
\varphi(\rho_1(u+v),\ldots,\rho_n(u+v))\leq
\varphi(\rho_1(u)+\rho_1(v),\ldots,\rho_n(u)+\rho_n(v))\tag*{(41)}
\]

En effet, c'est un cas particulier du Th.~1, $1^{\mathrm o}$, \unsure{hypothèse|hyp.} c, compte tenu des formules du N\textsuperscript{o}~1.

\deleted{\uline{Corollaire} \unsure{Soit|Soient} $u$ et $v$ étant comme ci-dessus, soit $\varphi(s)$ une fonction convexe \inserted{croissante} de $s\geq 0$, alors on a pour tout entier $n\geq 0$}

\[
\deleted{\sum_{1}^{n}\varphi(\rho_i(u+v))\leq\sum_{1}^{n}\varphi(\rho_i(u)+\rho_i(v))}\tag*{\deleted{(42)}}
\]

Nous serons surtout \unsure{intéressé|intéressés} par le

\uline{Corollaire} 1 Soient \deleted{$u$ et $v$ comme ci-dessus,}\inserted{$E$ et $F$ deux espaces de Hilbert,} et soit $N$ \uline{une norme de Schatten} sur $R^n$, \deleted{alors la quantité} posons, pour \deleted{tout} opérateur compact $u$ de $E$ dans $F$

\[
\lVert u\rVert_N=N(\rho_1(u),\ldots\rho_n(u))\tag*{(4\deleted{\unsure{2}}3)}
\]

Alors on a

\[
\deleted{\lVert\lambda u\rVert_N=|\lambda|\lVert u\rVert_N}\quad
\boxed{\lVert u+v\rVert_N\leq\lVert u\rVert_N+\lVert v\rVert_N}\tag*{(44)}
\]

\deleted{\unsure{et} $\unsure{u\to\lVert u\rVert_N}$ \unsure{définit|est donc} \illegible{2} $u\in L_0(E,F)$ \unsure{tels que|sous-espace} $N$ est une norme de Schatten \unsure{sur} $\unsure{R_+^{(n)}|R^{(N)}}$}

En effet, $N(s_1,\ldots s_n)$ pour $s_i\geq 0$ satisfait aux conditions de (41), on a donc

\[
N(\rho_1(u+v),\ldots,\rho_{\unsure{i|n}}(u+v))\leq
N(\rho_1(u)+\rho_1(v),\ldots,\rho_{\unsure{i|n}}(u)+\rho_{\unsure{i|n}}(v))
\]

et le second mb est majoré par $N(\rho_1(u),\ldots,\rho_n(u))+N(\rho_1(v),\ldots,\rho_n(v))$ puisque $N$ est une norme, d'où (44).

En particulier

\uline{Corollaire 1} \deleted{\unsure{Pour les|Soient}} si $u,v\in L_0^p(E,F)$, \deleted{\unsure{et si}} (cf $1\leq p\leq+\infty$) \unsure{on a} pour \deleted{\unsure{tout}} \inserted{entier $n>0$}

\[
\deleted{\lVert u+v\rVert_p\leq}
\left(\sum_{1}^{n}(\rho_i(u+v))^p\right)^{1/p}\leq
\left(\sum_{1}^{n}\rho_i(u)^p\right)^{1/p}+
\left(\sum_{1}^{n}\rho_i(v)^p\right)^{1/p}\tag*{(45)}
\]

donc \deleted{\unsure{en passant|par passage} \illegible{3}}

\[
\boxed{\lVert u+v\rVert_p\leq\lVert u\rVert_p+\lVert v\rVert_p}\tag*{(46)}
\]
% archive sheet number: 94
\scanpage{94}

\textcircled{5.} \uline{Classes remarquables d'opérateurs dans l'espace de Hilbert}

Nous commençons par donner un exposé de la théorie de Schatten, \deleted{\unsure{qui est}} sous une forme simplifiée rendue possible par la \deleted{\unsure{lemme}}\inserted{méthode du \unsure{minimum|minimax|convexité}} de Weyl, et en y apportant quelques compléments.

\margin{?}

\deleted{\unsure{Déf}} Rappelons la

\uline{Définition 1} \deleted{Soit} Une norme $N$ sur $R^n$ est dite \uline{norme de Schatten} si elle est symétrique, satisfait à \deleted{\unsure{à}}

\[
N(x_1,\ldots x_n)=N(|x_1|,\ldots |x_n|)\tag*{(1)}
\]

\deleted{et si elle} si elle est croissante \deleted{\unsure{sur l'ensem}} sur l'ens. des $x_i\geq 0$, et \deleted{\unsure{est satisf|elle satisfait}} est normalisée par $N(1,0,\ldots 0)=1$

Une norme $N$ sur $R^{(\underline N)}$ \deleted{\unsure{$\underline N$ ensemble}} (\unsure{espace|anneau} des suites infinies \deleted{\unsure{finies}} nulles sauf un nb fini de termes) est appellée une norme de Schatten, si \inserted{pour tout entier $n>0$,} elle induit \deleted{sur} \deleted{de} sur $R^n$ une norme de Schatten

($R^n$ sera ici identifié \deleted{à l'espace des suites}\inserted{au sous-espace} de $R^{(\underline N)}$ formé des suites nulles à partir du rang $n+1$). Notons qu'une norme de Schatten sur $R^n$ induit sur $R^m$ ($m\leq n$) une norme de Schatten. \deleted{\unsure{Tout cette théorie résulte que}}\inserted{Exemples de normes de} Schatten\,: \deleted{\unsure{Supposons}}

\[
\deleted{N_p(\{x_i\})=\left(\sum |x_i|^p\right)^{1/p}\quad\text{où }1\leq p}
\]

les normes $N_p((x_i))=\lVert(x_i)\rVert_p$, définies par \deleted{$1\leq p\leq\infty$ \unsure{par}}

\[
\lVert(x_i)\rVert_p=\left(\sum |x_i|^p\right)^{1/p}\qquad (1\leq p<+\infty)\tag*{(2)}
\]

\[
\lVert(x_i)\rVert_\infty=\underset{i}{\operatorname{Sup}}\,|x_i|\tag*{\deleted{\unsure{4}} (3)}
\]

Si $N_0$ est une norme de Schatten \deleted{alors} sur $R^n$, il existe toujours une norme de Schatten $N$ sur $R^{(\underline N)}$ \deleted{\unsure{qui}} induise $N_0$\deleted{, \unsure{p. ex.,} \illegible{1}} si, pour \deleted{\unsure{toute|tout}} $x=(x_i)\in\underline{c}_0$, on désigne par $(\unsure{\rho|f|R}_i(x))$ la suite des $|x_i|$ \deleted{\unsure{rangés}}\inserted{mais} rangés par ordre de grandeur décroissant, on peut \deleted{\unsure{p.ex.}} poser $N(x)=N_0(\rho_1(x),\ldots\rho_n(x))$. Il est trivial

\margin{\unsure{?}}
% archive sheet number: 95
\scanpage{95}

que $N(x)$ induit sur $R^n$ la fonction $N_0(x)$, et qu'on a $N(\lambda x)=|\lambda|N(x)$ et $N(x)=0\Rightarrow x=0$, il faut seulement prouver $N(x+y)\leq N(x)+N(y)$ (\unsure{en\textsuperscript{t}|est|seule} valable $=$ pour $x,y\in c_0$). Pour ceci, on \deleted{\unsure{regarde $(x_i)$ et $(y_i)$ comme les}}\inserted{remarque que si} $u_x$ est l'opération de multiplication par $x$ dans $\ell^2$ (qui est un opérateur cpt pour $x\in c_0$), on a \deleted{$u_{x+y}=u_x$} $\rho_i(u_x)=\rho_i(x)$, d'autre part \deleted{\unsure{que}} $u_{x+y}=u_x+u_y$, et l'inégalité à prouver résulte alors du N\textsuperscript{o}~4, th.~4, corollaire. En $=$ \deleted{\unsure{temps}}, les normes \inserted{de Schatten} ainsi associées à des normes \deleted{\unsure{de}} de Schatten $N_0$ sur $R^{n}$ \deleted{donnent} démontrent qu'il existe une infinité \inserted{continue} de normes de Schatten sur $R^{(\underline N)}$ équivalentes à la norme induite par $c_0$.

Soit $N$ \deleted{est} une norme de Schatten sur \deleted{\unsure{$R^n$} \unsure{finies}} $R^{(\underline N)}$. \deleted{Soit $\ell^N$ \unsure{l'espace des}}\inserted{Pour \deleted{toute}} suite \deleted{\unsure{réelle}} complexe $x=(x_i)$ \deleted{telle que \unsure{la suit} \unsure{on pose} soit $\unsure{N(x\ldots x)|N(x)}$}

\[
\deleted{N(|x_1|,\ldots |x_n|,0,0,\ldots)\quad\text{croît avec }n,\quad\text{on pose.}}
\]

Soit $|x|_n=(|x_1|,\ldots |x_n|,0,\ldots 0)$. \deleted{On a \unsure{aussitôt}}\inserted{Si $n$ varie} \deleted{alors $|x+y|_n\leq |x|_n+|y|_n$} $|x|_n$ est une suite croissante d'éléments positifs de $R^{(N)}$, posons

\[
N(x)=\lim_n N(|x|_n)\tag*{(4)}
\]

c'est un nb $\geq 0$ fini ou infini. La valeur \deleted{ainsi \unsure{désignée}}\inserted{ainsi} de $N(x)$ coïncide avec \deleted{\unsure{celle|$N(x)$}} \inserted{quand $x\in R^{(N)}$.} On a

\[
N(\lambda x)=|\lambda|N(x),\qquad N(\unsure{u|x}+y)\leq N(\unsure{u|x})+N(y)\tag*{(5)}
\]

\deleted{si $N(x)<+\infty$} \inserted{et dans} la première formule, \deleted{\unsure{il est à dire que son deuxième mb a un sens, i.e. $|\lambda|$}} \inserted{il faut prendre au deuxième} mb la valeur \uline{zéro} s'il \deleted{\unsure{est sous}}\inserted{se présente sous la} forme indéterminée $0\times\infty$. La deuxième formule résulte de ce que \deleted{\unsure{$N(x$}} $|x+y|_n\leq |x|_n+|y|_n$, d'où $N(|x+y|_n)\leq N(|x|_n+|y|_n)\leq N(|x|_n)+N(|y|_n)$.

\margin{\unsure{?}}

\deleted{Soit} Soit $\ell^N$ l'espace des suites complexes $x$ telles que $N(x)<\infty$, \unsure{muni évidemment des notions|espace vectoriel au moins|muni de} \deleted{\unsure{celles}} \unsure{des suites complexes nulles|de suite complexe nulle}, \unsure{partie|d'opération} \deleted{\unsure{de mult}} d'un \unsure{espace vect.|certain rang}. \unsure{Par|Sur|Pour} $(5)$, on voit que $N(x)$ est une \uline{norme} sur $\ell^N$.

\deleted{\unsure{De plus|D'ailleurs}, $\ell^N$ est un espace de \illegible{2} \unsure{positivité}. D'où \unsure{les inj} $R^{(N)}\in\ell^N$, et $N(x)$ induit \unsure{par hypothèse}, \unsure{donc} ($=$ \illegible{1} $\unsure{\alpha^{(i)}}$), est \illegible{2} \unsure{le candidat} dans $\ell^N$, \unsure{elles coïncident donc forcément}}
% archive sheet number: 96
\scanpage{96}

\deleted{On a, si $N$ \unsure{$<$|$\leq$}}

\[
\deleted{\ell^1\subset\ell^N\subset\ell^\infty}\tag*{\deleted{(6)}}
\]

\deleted{et la norme de $\ell^1$ est plus fine que celle induite par $\ell^N$, \unsure{qui est} et la norme de $\ell^N$ est moins fine que celle induite par}

Si $N_{\deleted{\unsure{1}}}'$ et $N_{\deleted{\unsure{2}}}''$ sont deux normes de Schatten telles que $N_{\deleted{1}}'\leq N_{\deleted{2}}''$, il est immédiat qu'on a $\ell^{N'_{\deleted{1}}}\subset\ell^{N''_{\deleted{2}}}$, et la norme de $\ell^{N'}$ \deleted{$\unsure{\ell^{N_1}\subset\ell^{N_2}|\ell^{N'}\subset\ell^{N''}}$} est plus fine que celle induite par $\ell^{N''}$.

\deleted{En particulier,} En particulier, si on remarque que toute norme de Schatten est comprise entre les normes de Schatten $N_1(x)=\sum_i |x_i|$ \deleted{\unsure{et}} $N_\infty(x)=\underset{i}{\operatorname{Sup}}\,|x_i|$, et que les espaces $\ell^N$ associés à ces dernières sont \deleted{$\unsure{\ell^1|\ell^{N_1}}$ et $\unsure{\ell^N|\ell^{N_\infty}}$} $\ell^1$ et $\ell^\infty$, \deleted{\unsure{est on}} voit que

\[
\ell^1\subset\ell^N\subset\ell^\infty\tag*{(6)}
\]

les applications identiques \deleted{\unsure{de $\ell^1$}} $\ell^1\to\ell^N$ et $\ell^N\to\ell^\infty$ étant de norme $\leq 1$.

\deleted{Notons la formule}

\[
\deleted{\sum\quad N\left(\sum_n x_{\unsure{i}}^{k}\right)\leq\sum_k N(x^k)}\tag*{\deleted{(7)}}
\]

\deleted{valable pour toute suite $\unsure{x^i=(x_i^i)|x^k}$ de suites $x^k\in C^N$, qui soit \deleted{\unsure{convergente}} sommable pour chaque coordonnée $i$\,: $\sum_k|x_i^k|<+\infty$. Il suffit de prouver (7) si le dernier mb est fini\,: alors, en posant $x=\sum_k x^k$, on aura $|x|_m\leq\sum_k|x^k|_m$ d'où $N(|x|_m)\leq N(\sum_k|x^k|_m)$}

% slanted
\margin{\unsure{Si Si}}

\deleted{or \inserted{$N$ étant une fonction continue sur $R^m$, on aura} $N(\sum_k|x^k|_m)=\lim_{\unsure{k|n}\to\infty}N(\sum_{k=1}^n|x^k|_m)$}

\deleted{et comme $N(\sum_{k=1}^{n}|x^k|_m)\leq\sum_{k=1}^n N(|x^k|_m)\leq\sum_{k=1}^{\unsure{\infty|n}} N(x^k)$, on aura $N(\sum_k|x^k|_m)\leq\sum_k N(x^k)$, donc pour tout $m$ on a $N(|x|_m)\leq\sum_k N(x^k)$, d'où par définition de $N(x)$ l'inégalité (7).}
% archive sheet number: 97
\scanpage{97}

On conclut de (7) que $\ell^N$ est \uline{complet}. Il suffit de prouver que \deleted{toute suite} $(x^k)$ \deleted{dans $\ell^N$} d'éléments $x^k\in\ell^N$, telle que $\sum\lVert x^{k}\rVert_N<+\infty$, est sommable dans $\ell^N$. Or \deleted{elle} on aura à fortiori $\sum\lVert x^k\rVert_\infty<+\infty$, donc la \deleted{série} suite sera sommable dans $\ell^\infty$, \inserted{vers un $x\in\ell^\infty$,} d'autre part l'on \deleted{aura} voit sur (7) que $x\in\ell^N$. D'autre part,

% slanted
\margin{Inutile, il suffit de procéder par dualité\,: $\ell^N$ est le dual de $\ell_N$.}

on a

\[
N\left(x-\sum_{k=1}^n x^k\right)=\deleted{\unsure{\leq}}N\left(\deleted{\unsure{x-|x}}\sum_{n+1}^\infty x^k\right)\leq\sum_{n+1}^{\infty}N(x^k),
\]

qui tend vers \deleted{\unsure{l'$\infty$} \unsure{puisque la|pour}}\inserted{$0$ pour $n\to+\infty$}, puisque $\sum N(x^k)$ est une série convergente. Il s'ensuit que l'on a aussi $\sum x^k=x$ \uline{dans $\ell^N$}.

Soit $\ell_N$ l'adhérence de $C^{(\underline N)}$ dans $\ell^N$, comme $\ell^N$ est complet, $\ell_N$ s'identifie au complété de $C^{(\underline N)}$ pour la norme $N$. Il est encore trivial que si $N'\leq N''$, alors $\ell_{N'}\subset\ell_{N''}$, et l'application identique $\ell_{N'}\to\ell_{N''}$ est de norme $\leq 1$. En particulier, les espaces $\ell_N$ associés à $N_1$ et $N_\infty$ étant resp. $\ell^1$ et $c_0$, on obtient

\[
\ell^1\subset\ell_N\subset c_0\tag*{(8)}
\]

\uline{Proposition} On a $\ell^N\subset c_0$ sauf si $N$ est équivalente à $N_\infty$:

\deleted{Notons la formule, qui précise la \unsure{condition|propriétés} \unsure{ci-dessus|suivantes}}

\[
\deleted{\underline{\ell}_N=\underline{\ell}^N\cap\underline{c}_0}\tag*{\deleted{(9)}}
\]

\deleted{Il suffit de prouver que $x\in\ell^N\cap c_0$ implique $x\in\ell_N$. \unsure{Et} En effet, on aura}

Notons que $\deleted{\unsure{C}}R^{(N)}$ est en dualité avec lui-$=$ par l'accouplement $\langle x,y\rangle=\sum x_i y_i$. Cela permet de faire correspondre à \deleted{toute} une norme $N$ sur $R^{(N)}$, la \uline{norme polaire}, notée $N'$, définie par

\[
N^\circ(y')=\underset{N(x)\leq 1}{\operatorname{Sup}}\,|\langle x,y'\rangle|\tag*{(9)}
\]

% margin note written sideways
\margin{De plus, \deleted{\unsure{comme|clairement}} comme $(N^\circ)^\circ=N$, on peut le vérifier .... \deleted{\unsure{on est ramené} \illegible{2}}}

Si $N$ est une norme de Schatten, on vérifie trivialement que $N^\circ$ est une norme de Schatten, \deleted{et} \unsure{On dira que} \deleted{le dual de $\ell_N$ s'identifie à $\ell^{N'}$, l'accouplement étant $\sum x_i x_i'$ est \unsure{abst} sommable, et que} \inserted{si $u\in\ell_N$, $u'\in\ell^{N^\circ}$, alors \deleted{\unsure{on a le}}}
% archive sheet number: 98
\scanpage{98}

\[
\sum|x_i x_i'|\leq N(x)N'(x')\unsure{)}\tag*{(10)}
\]

En effet, on a $\sum_{1}^{n}|x_i x_i'|=\langle |x|_n,|x'|_n\rangle\leq N(|x|_n)N(|x'|_n)\leq N(x)N'(x')$\unsure{)}, d'où (10). Il en résulte un \uline{accouplement} entre $\ell^N$ et $\ell^{N'}$, \deleted{\unsure{qui est}} d'où en restreignant \uline{entre $\ell_N$ et $\ell^{N'}$}\unsure{;|,} \deleted{\unsure{cet accoupl} on a ainsi associé à \unsure{$\ell$} de $\ell^{N'}$ une forme linéaire continue \unsure{de norme} $\leq 1$, $\ell^{N'}$ dans $(\ell_N)'$} \unsure{je dis} que \uline{le dual de $\ell_N$ s'identifie exactement à $\ell^{N'}$} (avec sa norme) par cet accouplement.

En effet, \deleted{\unsure{soit la}} il suffit de montrer qu'une forme linéaire continue $X'$ sur $\ell_N$ est définie par un \inserted{$x'$} élément de $\ell^{N'}$ \deleted{de \illegible{1}} tel que $N'(x')\leq\lVert X'\rVert$.

En effet, $X'$ \deleted{\unsure{est entièrement} \unsure{dé|aussi}} il suffit de le voir pour la \deleted{\unsure{part}} restriction à $\ell^1$, alors $X'$ \inserted{est} défini \deleted{\unsure{p|par}} \inserted{par} un $x'\in\ell^\infty$, \deleted{\unsure{avec}} on a $N'(|x'|_n)\leq\lVert X'\rVert$, d'où aussitôt $N'(x')\leq\lVert X'\rVert$, et le résultat.

-- Notons qu'en général, le dual de $\ell^N$ n'est pas $\ell^{N'}$, \deleted{\unsure{car}} car $\ell_N$ ne sera pas dense dans $\ell^N$. Ex\,: $N=N_\infty$. Notons que

\uline{Prop} \uline{pour que $\ell_N$ soit réflexif, il faut et il suffit que $\ell^N$ resp$^t$ $\ell_{N'}$ resp$^t$ $\ell^{N'}$ le soit}. \inserted{\uline{Il faut et il suffit pour cela que l'on ait $\ell_N=\ell^N$ et $\ell_{N'}=\ell^{N'}$.}} \deleted{En effet,}

$\ell^N$ réflexif implique $\ell_N$ réflexif (sous-espace fermé) qui implique $\ell^{N'}$ réflexif (dual de $\ell_N$) qui \inserted{de même} implique $\ell_{N'}$ réflexif et lui-$=$ implique $\ell^N$ réflexif, ce qui achève la démonstration \deleted{de l'équivalence pour} circulaire de l'équivalence des 4 conditions. Alors $\ell_{N'}$, qui est un sous-espace \deleted{\unsure{fort}} \uline{total} \inserted{et fortement fermé} du dual $\ell^{N'}$ de $\ell_N$, est identique à $\ell^{N'}$, et pour la $=$ raison on aura $\ell_N=\ell^N$.

\margin{?}

\inserted{Réciproquement,} s'il en est ainsi, alors le dual de $\ell_N$ est $\ell^{N'}=\ell_{N'}$, et le dual de ce dernier est $\ell^N=\ell_N$, donc $\ell_N$ coïncide avec son bidual, et est donc réflexif.

\inserted{Soit $N$ une norme de Schatten sur $R^{(N)}$ \deleted{non équivalente à $N_\infty$}\inserted{non équivalente à $N_\infty$}.} Soient $E,F$ deux espaces de Hilbert, \deleted{\unsure{si|$N$}} on désigne par $L^N_{\deleted{\unsure{0}}}(E,F)$ le \deleted{\unsure{sous|l'espace}} sous-espace de $L_0(E,F)$ formé des $u$ tels que $(\rho_i(u))\in\ell^N$, i.e. tels que

\[
N(u)=N((\rho_i(u)))\unsure{\leq|<}+\infty\tag*{(11)}
\]
% archive sheet number: 99
\scanpage{99}

\deleted{Si $N$ est équivalente à la} \deleted{Sous ces conditions,}

Alors la fonction $u\to N(u)$ sur $L^N_{\deleted{\unsure{0}}}(E,F)$ est une \uline{norme}. En effet, on a $N(u+v)=\lim_{n\to\infty}N(|(\rho_i(u+v))|_n\deleted{\unsure{,\rho_2(u+v)}})$, or le second mb est majoré par $N(|(\rho_i(u))|_n\deleted{\unsure{,\rho_2(u)}})+N(|(\rho_i(v))|_n\deleted{\unsure{,\rho_2(v)}})$ (N\textsuperscript{o} \unsure{,} formule (44)), donc par $N(u)+N(v)$. On note aussi trivialement $N(\lambda u)=|\lambda|N(u)$, et $N(u)=0\Rightarrow u=0$.

% slanted
\margin{$N(u)=N(u^{\unsure{*|+}})$ (12)\par $\left.\begin{aligned}N(uv)&\leq\lVert u\rVert N(v)\\N(vu)&\leq\lVert u\rVert N(v)\end{aligned}\right\}$ (13)}

L'espace $L^N_{\deleted{\unsure{0}}}(E,F)$ est donc un espace normé; \deleted{\unsure{prouvons qu'il est \uline{complet}.}}\inserted{soit} $L_N(E,F)$ le sous-espace normé formé des $u$ tels que $(\rho_i(u))\in\ell_N$. Je dis que $L_N(E,F)$ est l'adhérence de $\overline E\otimes F$ dans $L^N(E,F)$. En effet, montrons que

$\sum\rho_i\bar e_i\otimes f_i$ est limite au sens de $N$ de $\sum_1^n\rho_i\bar e_i\otimes f_i$, donc que $N(\sum_{n+1}^\infty\rho_i\bar e_i\otimes f_i)=N(\rho_{n+1}(u),\ldots)\to 0$ pour $n\to+\infty$, en d'autres termes, que dans $\ell_N$ les opérations $x\to t_n(x)=(\overbrace{0\ldots 0}^{n},x_{n+1},x_{n+2},\ldots)$ tendent vers $0$ simplement. Or elles forment une suite équicontinue, qui tend vers zéro dans \deleted{\unsure{un}} $R^{(N)}$, ce qui prouve la chose. \deleted{Notons \unsure{donc}, \unsure{lorsqu'on a} $L^N(E,F)$ \unsure{et que} \illegible{3}}

% margin note written sideways
\margin{Définition de $L^N(E,F)$ si $N\unsure{\simeq|=}N_\infty$ ($L^N(E,F)=L(E,F)$, avec \textcircled{14} $\boxed{N(u)=\underset{\substack{v\in F'\otimes E\\N^\circ(v)\leq 1}}{\operatorname{Sup}}\unsure{\operatorname{Tr}|\operatorname{tr}}(vu)}$ $L_N(E,F)$ comme avant). Cas où $u$ est cpt.}

\deleted{Le dual de $L_N(E,F)$ s'identifie à un sous-espace du dual de $L_{N_1}(E,F)=L^1(E,F)$, donc de $L(F,E)$, supposons que $N$ n'est pas équivalent à $N_1$, donc $N'$ non équivalent à $N_\infty$. Je dis que}

\[
\deleted{(L_N(E,F))'=L^{N'}(F,E).}
\]

\deleted{En premier lieu, \unsure{si} si $u\in L_N(E,F)$, $v\in L^{N'}(F,G)$ alors $uv\in L^1(E,G)$, avec}

\[
\deleted{\lVert uv\rVert_1\leq N(u)N'(v).}
\]

\deleted{soit} Soient $N,N',N''$ des normes de Schatten telles que

\[
\boxed{N(xy)\leq N'(x)N''(y)}\tag*{(\unsure{15|13})}
\]

pour tout $x,y\in R^{(N)}$. Alors la $=$ formule sera valable pour $x\deleted{\unsure{,y|y}}\in\ell^{N'}$, $y\in\ell^{N''}$. Prouvons alors que $u\in L^{N'}(E,F)$, $v\in L^{N''}(F,G)$ implique $uv\in L^N(E,G)$, avec

\[
\boxed{N(uv)\leq N'(u)N''(v)}\tag*{(16)}
\]
% archive sheet number: 100
\scanpage{100}

Cela résulte en effet \inserted{dans le cas où $N'\not\simeq N_\infty$, $N''\not\simeq N_\infty$,} du deuxième th de convexité pour opérateurs, et des formules $\rho_1(uv)\ldots\rho_n(uv)\leq(\rho_1(u)\rho_1(v))\ldots(\rho_n(u)\rho_n(v))$ (N\textsuperscript{o} \unsure{,} formule (37)). -- En particulier, \deleted{si $N=$ \unsure{la}} pour \deleted{\unsure{N}} deux normes polaires $N$ et $N^\circ$, on trouve

% slanted
\margin{Cas $N'\simeq N_\infty$ ou $N''\simeq N_\infty$\,: si l'on \unsure{adm.} ceci $u,v$ \uline{cpt} on s'y ramène facilement dans le cas général, car en tous cas $N(u)=\lim_P N(\unsure{P_n u|Pu})$ $P$ projecteurs orthogonaux de rang fini.}

\[
\boxed{\lVert uv\rVert_1\leq N(u)N^\circ(v)}\tag*{\textcircled{17}}
\]

On trouve aussi \deleted{direct$^{\mathrm{mt}}$} \deleted{\unsure{On a obtenu aussi} $\boxed{N(uv)\leq\lVert u\rVert N(v),\quad N(uv)\leq\lVert v\rVert N(u)}$}

Le dual de $L_N(E,F)$ s'identifie à un sous-espace du dual de $L_{N_1}(E,F)=L^{\unsure{1}}(E,F)$, donc de $L(F,E)$. Ce sous-espace contient $L^{N^\circ}(F,E)$, et l'application $L^{N^\circ}(F,E)\to(L_N(E,F))'$ est un isomorphisme \uline{métrique}.

Montrons que c'est \uline{sur}. On peut supposer que $N\not\simeq N_1$, prouvons que \deleted{\unsure{si}} \inserted{si un} $\Phi\in(L_N(E,F))'$, définit \inserted{$v\in L(F,E)$, \uline{on} a $v\in L^{N^\circ}(F,E)$. Or la norme de} \deleted{un op. \uline{compact} de $F$ dans $E$. \unsure{Comme} $v^*v$ \unsure{la dual} utilise $|\operatorname{Tr}vu|\leq M N(u)$ pour $u\in E'\otimes F$, or $w=v^*v\in L(E,E)$ aura la propriété analogue, prouvons que $w$ est}

$v=U|v|$, il suffit de prouver que $|v|\in L^{N^\circ}(E)$\unsure{,|$F$|$E$} or $|v|=U'v$, donc a la \unsure{$\subset$|même} propriété analogue à $v$ ($|\operatorname{Tr}(vu)|\leq M N(u)$ pour $u\in E'\otimes F$, $M$ constante). On est ramené à $F=E$, $v\geq 0$.

$v$ est cpt\,: sinon il existerait un op. continu $w$ tel que \deleted{\unsure{$v$}} $wv=P$, projecteur hermitien de rang infini, \deleted{\unsure{on aurait la}} on aurait encore $w\in(L_N(E,E))'$, ce qui est impossible, \deleted{\unsure{on suppose un projecteur}} \inserted{(en prenant} une base orthonormale dans $P$, \deleted{\unsure{en aurait}}\unsure{on voit qu'on|on soit qu'on} aurait \deleted{$\operatorname{Tr}\unsure{p_n}P=$} $1\in(\ell_N)'=\ell^{N^\circ}$, ce qui n'est pas. Donc

on aura $v=\sum\rho_i\bar e_i\otimes e_i$, et sur \uline{la \unsure{base|même|dite|trouve} formule} on voit que $(\rho_i)\in(\ell_N)'=\ell^{N^\circ}$, d'où la conclusion. -- Comme conséquence, on voit donc que $L^{N^\circ}(E,F)$ est \uline{toujours} complet, donc tout espace $L^N(E,F)$ est complet donc aussi $L_N(E,F)$

\margin{?}
% archive sheet number: 101
\scanpage{101}

\uline{Corollaire} 1 $u\in L^N(E,F)$ si et seulement si $|u|\in L^N(E,F)$, et

\[
N(u)=N(|u|)
\]

$u\in L_N(E,F)$ si et seulement si $|u|\in L_N(E,F)$.

\deleted{\uline{Corollaire 2} $N(u)$ est fonction croissante sur l'ensemble des op. hermitiens positifs dans $E$. En effet, $(\unsure{\rho|f}_i(u))$ est fonction croissante de $u$.}

\uline{Corollaire 2} Si $u=h+ik$, $h$ et $k$ hermitiens, alors

\[
u\in L_N(E)\Longleftrightarrow h,k\in L_N(E),\qquad
u\in L^N(E)\Longleftrightarrow h,k\in L^N(E).
\]

\uline{Remarque} Soit $p(u)$ une norme sur $E'\otimes F$ satisfaisant aux conditions\,: a)

\[
\begin{aligned}
p(Vu)&=p(u)\quad\text{pour $V$ unitaire dans $F$}\\
p(uU)&=p(u)\quad\text{--- $U$ --- $E$}
\end{aligned}
\]

b) $p(x'\otimes y)=p(x')p(y)$. Supposons $E$ et $F$ de dim infinie pour fixer les idées, \deleted{\unsure{soient}} $(e_i)$ $(f_i)$ des systèmes orthonormaux infinis dans $E$ et $F$, on a une immersion $R^{(\underline N)}\to E'\otimes F\subset L(E,F)$ par $(\lambda_i)\to\sum\lambda_i\bar e_i\otimes f_i$. Alors $p(\varphi(\lambda))$ sur $R^{(N)}$

est une norme, induisant sur $R^n$ des normes \deleted{de Schatten} invariantes par $\unsure{\mathfrak S_n|\mathscr G_n}$ (condition a sur $p$) normalisée (condition b) ne dépendant que de $(|\lambda_i|)$ (condition a). \deleted{\unsure{On a donc|Elle}} $N$ est donc une norme de Schatten. On vérifie alors immédiatement (condition a)) que l'on a $N(u)=p(u)$ pour \deleted{\uline{tout}} $u\in E'\otimes F$.

On a ainsi obtenu une caractérisation axiomatique des normes $N(u)$ sur les espaces \deleted{$L(E,F)$.} $E'\otimes F$.
% archive sheet number: 102
\scanpage{102}

\fbox{Espaces de Schatten}

\textcircled{1} \uline{Normes de Schatten}

\uline{Définition}. \uline{Prop} Une norme de Schatten $N(x)$ $\unsure{\to|\text{est}}$ croissante pour $x$ positif. \unsure{\}|§|?}

\uline{Exemples} les $N_p$. Norme de Schatten définie par une Norme de S. sur $R^n$.

\uline{Définition} $N(x)$ pour $x\in\underline C^{\underline N}$. $\ell^N\ni x\Longleftrightarrow N(x)$ fini, $N$ est une norme sur $\ell^N$. Espace $\ell_N$, adhérence de $C^{(\underline N)}$ dans $\ell^N$

\uline{Inclusions} pour $N'\leq N''$, en particulier

\[
\ell^1\subset\ell^N\subset\ell^\infty,\qquad \ell^1\subset\ell_N\subset c_0
\]

1 \uline{Prop} $\ell_N\subset\underline c_0$ sauf si $N$ est équivalente à $N_\infty$

\uline{Polarité}, $(N^\circ)^\circ=N^\circ$\,: \uline{Prop} le dual de $\ell_N$ est $\ell^{N^\circ}$. \uline{Corollaire} $\ell^N$ est complet, $\ell_N$ est complet. \uline{Prop} $\ell_N$ réflexif $\Longleftrightarrow\ell^N$ réflexif $\Longleftrightarrow\ell_{N^\circ}$ réflexif $\Longleftrightarrow\ell^{N^\circ}$ réflexif $\Longleftrightarrow\ell_N=\ell^N$ et $\ell_{N^\circ}=\ell^{N^\circ}$.

\deleted{\unsure{Exemples}} \uline{Cas des $N_p$} On a $(N_p)^\circ=N_{\!p'}^{\deleted{\unsure{p}}}$. On retrouve la réflexivité des $\ell^p$ $(1<p<+\infty\unsure{)}$

\textcircled{2} \uline{Espaces de Schatten}

\deleted{\uline{Définition} Si $u\in L_0^{\unsure{N}}(E,F)$, \inserted{$N(u)=$}$N((\rho_i(u)))$\unsure{;} $u\in L^N(E,F)\Longleftrightarrow N(u)<+\infty$, \unsure{Prop} $N(u)$ est une norme sur $L^N(E,F)$, $L_N(E,F)$ adhérence de $E'\otimes F$ ($N$ non équivalente à $N_\infty$). $N(u)=N(u^*)$, $N(uv)\leq\lVert v\rVert N(u)$, $N(vu)\leq\lVert v\rVert N(u)$.}

\deleted{\uline{Définition} (Si $N$ équivalente à $N_\infty$) $N(u)=\underset{\substack{N^\circ(v)\leq 1\\v\in F'\otimes E}}{\operatorname{Sup}}|\operatorname{Tr}vu|$, Si $u\in L_0(E,F)$, on retrouve $N(u)=N((\rho_i(u)))$, i.e, pour tout $N$}

\[
\deleted{\unsure{(}N((\rho_i(u)))=\underset{\substack{v\in F'\otimes E\\N^\circ(v)\leq 1}}{\operatorname{Sup}}|\operatorname{Tr}vu|}
\]

\uline{Définition\,:} Si $u\in L_0(E,F)$, posons $N(u)=N(\rho_i(u))$

\[
N(u)=\deleted{N((\rho_i(u)))}=\underset{\substack{v\in F'\otimes E\\\deleted{\unsure{N^\circ((\rho_i(v)))\leq 1|N^\circ(v)\leq1}}\\N^\circ(v)}}{\operatorname{Sup}}|\operatorname{Tr}vu|
\]

Définissons, pour $u\in L(E,F)$, $N(u)$ par la formule ci-dessus. Soit $L^N(E,F)\ni u\Longleftrightarrow\deleted{\unsure{N(u)=\operatorname{Sup}\operatorname{Tr}vu}}N(u)<+\infty$, $N\inserted{(u)}$ \unsure{est} une \uline{norme} sur $L^N(E,F)$ \deleted{et on a} soit $L_N(E,F)$ l'adhérence de $E'\otimes F$ dans $L^N(E,F)$

\deleted{identique aux} \uline{Prop} $u\in L_N(E,F)\Longleftrightarrow u\in L_0(E,F)$ et $(\rho_i(u))\in\ell_N$ \deleted{\unsure{Cor}} \deleted{\uline{Prop}} \uline{Inclusions} pour $N'\leq N''\Rightarrow\deleted{\unsure{\ell^{N''}\subset\ell^{N'}}}\,L^{N''}(E,F)\subset L^{N'}(E,F)$ et $L_{N''}(E,F)\subset L_{N'}(E,F)$, en particulier

\[
\begin{aligned}
L^1(E,F)&\subset L^N(E,F)\subset L(E,F)\\
L^1(E,F)&\subset L_N(E,F)\subset L_0(E,F)
\end{aligned}
\]

\deleted{\uline{Prop} Si $N$ non équivalente à $N_\infty$, alors}

\[
\text{Prop}\quad\boxed{N(u)=N(u^*)\quad N(uv)\leq\lVert v\rVert N(u),\ N(vu)\leq\lVert v\rVert N(u)}
\]

\deleted{\uline{Caractérisation} \unsure{$N(u)$, utilisant.}} (par transposition il suffit de le voir pour $u\in L_0(E,F)$, celà résulte de $\rho_i(vu)\leq\lVert v\rVert\rho_i(u)$, \ldots)
% archive sheet number: 103
\scanpage{103}

\uline{Corollaire} $N(u)=N(|u|)$\,: si $u=h+ik$, $N(h),N(k)\leq N(u)\leq N(h)+N(k)$, \unsure{obtient|donc} $u\in L^N(E)$ \deleted{$\Longleftrightarrow$}$\Longleftrightarrow h,k\in L^N(E)\deleted{\unsure{(F)}}$. \deleted{\unsure{et}}

% slanted
\margin{Analogue pour $L_N(E,F)$}

\uline{Prop} $L^N(E,F)\subset L_0(E,F)$ si $N$ non équivalente à $N_\infty$\unsure{;} (sinon $L^N(E,F)=L(E,F)$). En effet, si $u\in L^N(E,F)$, \deleted{prouvons que}\inserted{alors} $|u|\in L^N(E)$, prouvons $|u|$ compact (d'où $u$ compact). Sinon, il existerait un op. hermitien $v\in L(E)$ tel que $v|u|$ soit un projecteur hermitien infini, or on aurait $v|u|\in L^N(E)$. C'est impossible, car on \deleted{\unsure{aura $(\unsure{\rho|f}_i)\in\ell^N$|aurait $(\rho_i)\in\ell^N$}} aurait $\sum|x_i|\deleted{\unsure{1}}\leq M$ pour $N^\circ((x_i))\leq 1$, donc $N^\circ$ équivalente à $N_1$, ce qui est absurde.

\uline{Corollaire} $L^{N_p}(E,F)=L^p(E,F)$ pour $1\leq p\leq+\infty$\,: \deleted{si $p<+\infty$, est}\inserted{si $p<+\infty$, est} même identique à $L_{N_p}(E,F)$. (car $\ell_N=\ell^N=L_N(E,F)=L^N(E,F)$\unsure{)}

\uline{Prop} $u\in L^N(E,F)$, $v\in L^{N^\circ}(F,G)\Rightarrow vu\in L^1(E,G)$

D'où accouplement Trace entre $L^N(E,F)$ et $L^{N^\circ}(F,E)$. Par restriction, \deleted{$L^{N^\circ}(F,E)$ est un sous-espace du dual} \unsure{en principe|on prouve}, on a

\uline{Prop} le dual de $L_N(E,F)$ est $L^N(E,F)$

\uline{Corollaire 1} $L^N(E,F)$ et $L_N(E,F)$ sont complets

\uline{Corollaire 2} $L_N(E,F)$ réflexif $\Leftarrow N$ \uline{réflexif}. Réciproque vraie si $E$ et $F$ de dimension infinie. \uline{En tous cas} $L_N(E,F)$ réfl. $\Longleftrightarrow L^N(E,F)$ réflexif $\Longleftrightarrow\ldots$

\uline{Corollaire 3} Si $1<p<+\infty$, alors $L^p(E,F)$ est un espace réflexif ayant pour dual $L^{p'}(F,E)$

\uline{Prop} Si $N,N',N''$ sont des normes de Schatten telles que $N(xy)\leq N'(x)N'(y)$, alors $u\in L^{N'}(E,F)$, $v\in L^{N''}(F,G)$ implique $vu\in L^N(E,F)$, \inserted{et $N(vu)\leq N'(u)N''(v)$ \uline{Corollaire}} Si on met \deleted{\unsure{soit}} $N',N''$ \unsure{est} mis en bas, on peut faire de $=$ avec l'autre.

(\deleted{par transposition, on se ramène \unsure{au deuxième} \illegible{2}, en} \deleted{prenant $N'''=N^\circ$,} Par le deuxième th. de convexité si $v$ et $u$ sont \deleted{\unsure{là par contraction}, on} \deleted{à prouver que si $N_1(xyz)\leq N'(x)N''(y)N'''(z)$, alors} cpts. On s'y ramène facilement, en notant que \deleted{on a $N_1(wvu)\leq N'(u)N''(v)N'''(w)$.}

\[
N(u)=\lim N(\unsure{pu|pup}\deleted{\unsure{p}})\unsure{p}=\lim_p N(pup)=\lim_p N(up)
\]

$p$ projecteur orth. de rang fini, variant en filtrant croissant. (Car $N(pvup)=N((pv)(up))\leq N''(pv)N^{\unsure{''|'}}(up)\leq N''(v)N^{\unsure{\circ|'}}(u)$). Enfin, rappelons

\uline{Prop} $N((\lambda_i(u)))\leq N(u)$

\uline{Prop} Caractérisation des espaces de Schatten.
% archive sheet number: 104
\scanpage{104}

\uline{Théorème} \textcircled{1} Soit $N$ une norme de Schatten sur $R^{(\underline N)}$, \deleted{Soient $E$ et $F$}\inserted{pour tout couple $(E,F)$ de} deux espaces de \deleted{Banach}\inserted{Hilbert}, \deleted{posons}\inserted{et} pour $u\in L_0(E,F)$, posons

\[
N(u)=N((\rho_i(u)))\tag*{(18)}
\]

et pour $u\in L(E,F)$

\[
N(u)=\underset{\substack{v\in F'\otimes E\\N^\circ(v)\leq 1}}{\operatorname{Sup}}|\operatorname{Tr}vu|\tag*{(19)}
\]

Cette \unsure{deuxième|dernière} expression coïncide avec la précédente sur $L_0(E,F)$. Soit $L^N(E,F)$ le sous-espace de $L(E,F)$ formé des $u$ tels que $N(u)<+\infty$, alors $N$ est une norme sur $L^N(E,F)$, qui en fait un espace de Banach. \deleted{L'adhérence \inserted{$L_N(E,F)$} de $E'\otimes F$ dans cet espace est l'espace des $u\in L_0(E,F)$ tels que} Si $u\in L_0(E,F)$, alors on a $u\in L^N(E,F)$ si et seulement si $(\rho_i(u))\in\ell^N$, et alors on a (18). \deleted{\unsure{N}} L'adhérence \inserted{$L_N(E,F)$} de $E'\otimes F$ dans $L^N(E,F)$ est l'ensemble des op. \deleted{\unsure{cpts}} $u\in L_0(E,F)$ tels que $(\rho_i(u))\in\ell_N$.

\textcircled{2} On a

\[
\boxed{N(u)=N(u^*)}\deleted{\unsure{N(u)}}\tag*{\textcircled{20}}
\]

et par suite $u\in L^N(E,F)\mathrel{\uline{=}}u^*\in L^N(F,E)$ \deleted{\unsure{$\Longleftrightarrow u\in L^N(E,F)$|sur $L^N(E,F)$}}\inserted{de $=$} $u\in \unsure{L_N|L_N^c}(E,F)\Longleftrightarrow u^*\in L_N(F,E)$. On a

\margin{?}

\[
\begin{aligned}
N(uv)&\leq\lVert u\rVert N(v)\\
N(\unsure{uv|vu})&\leq\lVert v\rVert N(u)
\end{aligned}\tag*{\textcircled{21}}
\]

si $u\in L(E,F)$, $v\in L(F,G)$. \deleted{\unsure{Donc}} En composant à gauche ou à droite un op. \deleted{$\in L^N($} du type $\ell^N$ (resp$^t$ $\ell_N$) avec un opérateur \deleted{de type} \uline{continu}, on obtient un op. du type $\ell^N$ (resp.t $\ell_N$). \inserted{$N(u)$ est fonction croissante sur l'ens. des hermitiens positifs dans $L_0(E)$.}

\textcircled{3} \deleted{\textcircled{4}} Si $N,N',N''$ sont \deleted{deux}\inserted{des} normes de Schatten telles que

\[
\boxed{N(xy)\leq N'(x)N'(y)}\tag*{\textcircled{22}}
\]

pour $x,y\in R^{(\underline N)}$, alors on a, pour $u\in L(E,F)$, $v\in L(F,G)$,

\[
N(vu)\leq N'(v)N''(u)\tag*{\textcircled{23}}
\]

donc si $u\in L^{N'}(E,F)$, $v\in L^{N''}(E,F)$, alors $vu\in L^N(E,F)$, et si on a \inserted{de plus} $u\in L_{N'}(E,F)$ ou $v\in L_{N''}^{\deleted{\unsure{N}}}(E,F)$, on aura $=$ $vu\in L_N(E,F)$, -- En particulier, on aura, pour $u\in L(E,F)$, $v\in L(F,G)$

\[
\lVert vu\rVert_1\leq N(u)N^\circ(v)\tag*{\textcircled{24}}
\]

donc si $u\in L^N(E,F)$, $v\in L^{N^\circ}(F,G)$, on a $vu\in L^1(E,G)$, et

\[
|\operatorname{Tr}vu|\leq N(u)N^\circ(v).\tag*{\textcircled{25}}
\]
% archive sheet number: 105
\scanpage{105}

\textcircled{4} Par l'accouplement \deleted{\unsure{de}}

\[
\langle u,v\rangle=\operatorname{Tr}vu\tag*{\textcircled{26}}
\]

le dual de $L_N(E,F)$ s'identifie à $L^{N^\circ}(F,E)$. Si $E$ et $F$ sont de dimension infinies, \deleted{\unsure{les espaces suivants} $L_N(E,F)$} \inserted{alors} les conditions suivantes sont équivalentes\,: a) $L_N(E,F)$ est réflexif b) $L^N(E,F)$ est réflexif \deleted{$L_N(E,F)=L^N(E,F)$ et $L_{N^\circ}(F,E)=L^{N^\circ}(\unsure{F|F,E})$} c) $\ell_{\deleted{\unsure{\ell}}N}$ est réflexif \deleted{\unsure{f),}} d) $\ell_{N'}$ est réflexif\unsure{.|;} \deleted{\unsure{c'est et seulement si}, \unsure{$\ell$} équiv$^{\mathrm{tes}}$ de ces propositions à\,: d) $\ell^N$ est réflexif. -- e) $\ell_{N^\circ}$ est réflexif et \unsure{enfin} f) $\ell^{N^\circ}$ est réflexif)} \deleted{alors} d) $L_N(E,F)=L^N(E,F)$ et $L_{N^\circ}(F,E)=L^{N^\circ}(F,E)$

\textcircled{5} \deleted{\unsure{Remarquons|Remarque}} Soit $u\in L_0(E,F)$, \deleted{si $u\in L^N(E,E)$, \unsure{alors} \unsure{Alors} $(\lambda_i(u))\in\ell^N$,} \inserted{Alors}

\[
N((\lambda_i(u)))\leq\deleted{N((\unsure{\rho|f}_i(u)))}N(u)\tag*{(27)}
\]

Si $u\in L^N(E,F)$, \deleted{alors (resp$^t$ $L_N(E,E)$)} alors, $(\lambda_i(u))\in\deleted{\unsure{L^N\!,|\ell_N}}\ell^N$ \deleted{(resp$^t$ $\ell_N$)}

Tout est démontré, sauf 4. \deleted{\unsure{et}} et 5. Pour 4, on note\,: si $E$ et $F$ sont de dimension infinie, alors $\ell^N$ \inserted{(resp.t $\ell_N$)} est isométrique à un sous-espace vectoriel normé de $L^N(E,F)$ (resp.t $L_N(E,F)$. (prendre $(e_i)$, $(f_i)$ systèmes orthonormaux infinis, et $(\lambda_i)\to\sum\lambda_i\bar e_i\otimes f_i$ de $\ell^N$ dans $L^N(E,F)$). \deleted{Donc}

Si $L^N_{\deleted{\unsure{0}}}(E,F)$ est réflexif, $L_N(E,F)$ est réflexif, et alors aussi $\ell_N$ est réflexif, \deleted{\unsure{donc}}\inserted{Celà implique} (voir --) $\ell_N=\ell^N$, $\ell_{N^\circ}=\ell^{N^\circ}$, donc $L^N(E,F)=L_N(E,F)$ (car on \deleted{\unsure{peut supposer}}\inserted{aura} \unsure{exigé|déjà} $N\not\simeq N_\infty$), donc $L^N(E,F)$ sera réflexif, d'où l'équivalence a), b), c).

Et le fait que cela implique d) (d'abord la première \unsure{relation|assertion}, pour la seconde \unsure{on|car} $\ell_{N^\circ}$ sera aussi réflexif). Enfin d implique \deleted{$\ell_N=\ell^N$ et $\ell_{N^\circ}=\ell^{N^\circ}$ (voir \unsure{théorème précédent}) donc $\ell_N$ réflexif (toujours)} \inserted{que le dual de $L_N(E,F)$ est $L_{N^\circ}(F,E)$,} dont le dual est $L_N(E,F)$, donc $L_N(E,\unsure{F|R})$ est réflexif.

\textcircled{5} (27) Est le corollaire 5 du 1. th. de convexité \deleted{\unsure{Résulte aussitôt du|résulte immédiatement}}
% archive sheet number: 106
\scanpage{107}

Décomposition en facteurs des fonctions entières.

\uline{Théorème 1} (Weierstrass) Soit $f(z)$ fonction entière, $(a_n)$ la suite de ses zéros, rangés par ordre en modules croissants. Supposons $f(0)\ne0$. Il est possible de trouver une suite croissante de nombres entiers $\geq0$ $p_n$, et une fonction entière $g(z)$, tels que l'on ait

\[
 f(z)=e^{g(z)}\prod_n\left(1-\frac{z}{a_n}\right)e^{\frac{z}{a_n}+\frac{z^2}{2a_n^2}+\cdots+\frac{z^{p_n}}{p_n a_n^{p_n}}} \tag*{(1)}
\]

le produit étant uniformément absolument convergent quand $z$ parcourt n'importe quelle région bornée. Si \deleted{de plus} \inserted{une suite d'entiers} $p_n$ vérifie

\[
 \sum\left|\frac{r}{a_n}\right|^{p_n+1}<+\infty\qquad\text{quel que soit }r,
\]

alors on peut employer cette suite dans l'expression \uline{1}. En particulier, supposons $k$ un entier tel que $\sum\frac{1}{|a_n|^{k+1}}<+\infty$, alors on peut prendre $p_n\deleted{\unsure{=|k}}\unsure{=}p$ pour \deleted{tout} $n$, avec $p\leq k$ d'où

\[
 f(z)=e^{g(z)}\prod_n\left(1-\frac{z}{a_n}\right)e^{\frac{z}{a_n}+\cdots+\frac{z^k}{ka_n^k}} \tag*{(2)}
\]

% archive sheet number: 108
